\documentclass[12pt,a4]{amsart}


%\pagestyle{empty}
\usepackage[T1]{fontenc}
\usepackage{palatino}
\usepackage{mathpple}

\newcommand{\speaker}[1]{\par\noindent{\bf #1}\par\medskip}
\renewcommand{\title}[1]{\par\noindent{#1}\par\medskip}
\renewcommand{\abstract}[1]{\par\noindent{Abstract: }#1\par\medskip}
\begin{document}
\begin{center}
{\Large \bf Abstract}
\\
\medskip
\end{center}

\bigskip
%%%%%%%%%%%%%%%%%%%
\bigskip

\speaker{%
Eric Bedford (Indiana University, USA, and RIMS)}
\title{Complex Methods for Dynamical Systems}
\bigskip

\speaker{%
Tatsuya Tate (Keio University)}
\title{Distribution laws for integrable eigenfunctions}
\abstract{%
We give several precise results on the asymptotic behavior of
the distribution functions of the joint eigenfunctions of the completely
integrable quantum system defined by the Hamiltonian torus action
on a toric K\"ahler manifold.
We show that there is a universal Gaussian scaling limit
of the distribution functions near its center of the localization.
We also determine the limit distribution for the tail of the localization.
It is not a universal result but it depends on the global geometry
of the toric K\"ahler manifold, and in particular it depends on
the rate of the exponential decay of the eigenfunctions
away from the classically arrowed region.
}

\bigskip


\speaker{Miroslav Engli\v s (Academy of Sciences of the Czech Republic)}   

\title{A characterization of symmetric domains}

\abstract{%
Takaaki Nomura has recently proved that if the Berezin transform on a
   homogeneous complex domain commutes with the invariant Laplacian, then
   the domain must be symmetric. We show that a similar characterization
   of symmetric domains prevails even for the much wider class of all
   pseudoconvex domains with a K\"ahler metric: namely, if the Berezin
   transform commutes with the Laplace-Beltrami operator and at the same
   time is formally self-adjoint with respect to the K\"ahler measure, then
   the domain must be symmetric and the K\"ahler metric must be a constant
   multiple of the invariant (i.e. Bergman) metric.
}

\newpage

\speaker{Richard Penney (Purdue University, USA)}
\title{Harmonic functions (in the sense of Helgason) on bounded (not necessarily symmetric) homogeneous domains in $\mathbb{C}^n$}

\abstract{%
Let $X=G/K$ be a Riemannian symmetric space and let $D_G(X)$ be the algebra  of all
$G$-invariant differential operators on $X$.  Let $\chi$ be a character of $D_G(X)$.   A $C^\infty$ function $F$ on $X$ is {\it
$\chi$-harmonic} if
\begin{equation}
XF=\chi(X)F
\end{equation}
for all $X\in D_G(X)$.

One of the
most
beautiful results in the harmonic analysis of symmetric spaces is the so called ``Helgason
Conjecture,'' which was finally proved in full generality by M. Kashiwara, A. Kowata, K. Minemura,
K. Okamoto, T. Oshima, and M. Tanaka  which states that on a Riemannian symmetric space of
non-compact
type,  a function satisfies (1) if and only if it is the Poisson
integral of a hyperfunction over the Furstenberg boundary. Furthermore, a result of due to Oshima and Sekiguchi shows that the boundary values of $F$ are distributions if and only if $F$ satisfies
\begin{equation}
|F(x)|\le Ae^{r_o\tau(x)}
\end{equation}
for all $x\in X$ where $\tau(x)$ is the Riemannian distance in $X$ from
$x$ to the
base point $x_o=eK$.

In this talk we discuss a generalization of this result to bounded homogeneous domains in $\mathbb{C}^n$ in which the algebra $D_G(X)$ is replaced by an algebra $D_{geo}(X)$  of $G$-invariant differential operators canonically defined from geometry of $X$. We prove that functions which satisfy (1) and (2) have distributional boundary values and that there is a well defined ``Poisson Transformation'' which produces $F$.  Our results are based on producing generalizations of the asymptotic expansions studied by N. Wallach and by E. van den Ban and H. Schlichtkrull.
}
\newpage



\speaker{%
Hideyuki Ishi (Yokohama City University)}
\title{
Matrix realization of bounded homogeneous domains}
\abstract{% 
The Siegel disk is a complex bounded domain in the space of 
complex symmetric matrices. It is a natural generalization of 
the unit disk, and besides, we shall show that it has a universal
nature in a sense among homogeneous bounded domains. Namely,
based on the one-to-one correspondence between homogeneous 
bounded domains and normal j-algebras, we construct 
a canonical homogeneous imbedding of every homogeneous bounded domain 
into a direct product of the Siegel disks. Then we describe the
Bergman metric on the homogeneous bounded domain explicitely
in terms of the determinant and the trace of matrices.
}
\bigskip


\speaker{%
Genkai Zhang (Chalmers University of Technology and University of Gothenbur, Sweden) }
\title{Invariant plurisubharmonic
functions on extended Cartan domains
}
\abstract{\def\GG{\mathbb G}
\def\KK{\mathbb K}
Let $D=\GG/\KK$ be an irreducible bounded
symmetric tube domain in a complex $d$-dimensional vector space $V$
and let $\mathcal S$
be its Siegel domain realization. Let $K$
be the  semi-simple factor of $\KK$. The group $K$
and its complexification
$K_{\mathbb C}$ act linearly on $V$.
We introduce the extended Cartan domain
$D_N^\prime$
 as the orbit of the diagonal action of $K_{\mathbb C}$
on the product $D^N$, $D_N^\prime=K_{\mathbb C}\cdot D^N$. We prove
that $D_N^\prime$ is a domain of
holomorphy for all rank two domains $D$ and for $N>d$.
We also consider similar problems
for the extended Siegel domains.
}


\bigskip


\speaker{%
Mihai Paun (Universit\`e Louis Pasteur, France)}
\title{Numerical properties of currents associated 
to transcendental entire curves}
\abstract{Consider a compact K\"ahler 
manifold
$X$, and let $\varphi:\mathbb{C}\mapsto X$ be a 
Zariski dense entire curve. We denote by 
$T[\varphi]$ the associated current (cf.\
McQuillan, Publ.\ IHES, 87, 1998). We show that
the current $T[\varphi]$ is semi-positive on 
the nef cone of $X$. If $T[\varphi]$ acquires 
a special kind of singularity (called 
multiplicity) at some point of the manifold 
$X$, then we show that $T[\varphi]$ is 
strictly positive on the nef cone. }

\newpage

\speaker{%
S\'ebastien Boucksom (Universit\`e  Paris 7, France)}
\title{Characterisation of pseudo-effectivity via covering curves}

\abstract{%
A holomorphic line bundle over a compact complex manifold is said to be
pseudo-effective if it can be endowed with a singular hermitian metric whose
curvature current is positive. In a recent joint work with Demailly, Paun and
Peternell, we have given the following algebraic characterisation of
pseudo-effectivity: a line bundle over a projective manifold is pseudo-effective
iff it has non-negative degree over every covering family of curves. The proof
is based on the construction of ``moving intersection classes" of
pseudo-effective cohomology $(1,1)$-classes, which are morally the intersections
of their nef parts in would-be Zariski decompositions, and the generalisation to
that setting of the following well-known facts in the 2-dimensional situation:
the Zariski decomposition $D=P+N$ of a pseudo-effective divisor $D$ into a nef part
$P$ and a contractile divisor $N$ exists, is orthogonal, and the self-intersection
of $P$ computes the volume of $D$. The technical setting is that of currents with
minimal singularities and their current-theoretic intersection. 
}
\bigskip

\speaker{%
Yoshinori Namikawa (Kyoto University)}
\title{%
Counter-examples to global Torelli 
problem for irreducible symplectic 
manifolds}

\bigskip

\speaker{Sergey Pinchuk (Indiana University, USA)}

 \title{Regularity of CR mappings}

\abstract{%
I will discuss the main ideas of the proof of the following joint result with Klas Diederich.

\medskip

\noindent
{THEOREM. }{\em  Let $M, M'$ be smooth real analytic hypersurfaces in ${\mathbb C}^n$ of finite type 
(in the sense of D'Angelo) and  $f:M\to M'$ be a continuous CR mapping. Then
$f$ extends holomorphically to a neighborhood of $M$.
}


\medskip

This result is new even under additional assumptions that $M, M'$ are pseudoconvex
and $f$ is smooth. The main tools of the proof are the techniques of Segre varieties and 
convergence properties of analytic sets. In particular, some uniform volume 
estimates are established for holomorphic families of analytic sets.
}


\newpage

\speaker{%
Akio Kodama (Kanazawa University) and Satoru Shimizu (Tohoku University)}
\title{A group-theoretic characterization
of the space obtained by omitting the coordinate hyperplanes from
the complex Euclidean space}
\abstract{In this talk, we discuss a characterization of
$\mathbb{C}^k \times(\mathbb{C}^{\ast})^{\ell}$
by its holomorphic automorphism group.
We show that if a connected holomorphically separable
complex manifold $M$ of dimension $n$
has the holomorphic automorphism group $\text{\rm Aut}(M)$
isomorphic to
$\text{\rm Aut}
 (\mathbb{C}^k
 \times
 (\mathbb{C}^{\ast})^{n - k})$
as topological groups, then $M$ itself is biholomorphically
equivalent to
$\mathbb{C}^k
 \times
 (\mathbb{C}^{\ast})^{n - k}$.
We also give a new approach to the study of $U(n)$-actions on
complex manifolds of dimension $n$.
}

\bigskip


\speaker{%
Evgeni Chirka (Steklov Mathematical Institute of Russian Academy of Science)
}
\title{On the extension of holomorphic motions}
\nopagebreak
\abstract{%
Let $\{f_\tau\}_{\tau\in E}$ be a family of holomorphic functions with pairwise disjoint graphs on a complex manifold $\Omega$ and $\tau=f_\tau(z_0)\in E\subset \mathbb{C}$ for a base point $z_0\in\Omega$ (i.e., $f(z\tau)=f_\tau(z)$ is a {\em holomorphic motion} of the set $E\subset\mathbb{C}$ over $\Omega$). The problem is in extending of this family to the holomorphic motion of the whole complex plane $\mathbb{C}$.

For $\Omega=\mathbb{D}$, the unit disc in $\mathbb{C}$, the problem was solved by Slodkowski (1991) using the solvability of special nonlinear Riemann--Hilbert problem, so the specificity of the disc was used essentially.

In general case, there is natural topological obstruction for the extendability of some motions. If $\{f_\tau\}$ is normalized by the conditions $0,1\in E$, $f_0\equiv0, f_1\equiv 1$, the condition $\Delta_\gamma \arg f_\tau=0$, for an arbitrary closed curve $\gamma\subset\Omega$ and any $\tau\ne0$ in $E$, is necessary and sufficient for the continuous extension of the motion (i.e., with continuous $f_\tau\in\mathbb{C}$).

We show that, for an arbitrary {\sl Riemann surface} $\Omega$, this topological condition is also sufficient for the holomorphic extendability (thus, the Oka's principle take place for the extending of holomorphic motions over an arbitrary Riemann surface). The proof is reduced to the solvability of the special quasi-linear $\overline\partial$-equations on compact Riemann surfaces.

Some generalizations and open problems will be discussed as well.
}


\newpage


\speaker{%
Klas Diederich (University of Wuppertal, Germany)}
\title{Complex Analysis on lineally convex domains}
\abstract{%
One of the important goals of quantitative complex analysis
is to do $\overline\partial$-theory with $L^p$- and H\"older
estimates on a large class of domains in $\mathbb{C} ^n$.
The main method available for this consists in the construction
and estimation of suitable $\overline\partial$-solving integral
kernels based on the existence of suitable Leray sections for the
domains. The necessary smooth families of holomorphic support
functions have been constructed by Diederich-Fornaess for all
lineally convex bounded domains in $\mathbb{C} ^n$ with smooth
boundaries of finite type. We will consider in the lecture this
class of domains and show how it is possible to construct a suitable
pseudometric on them with the help of which we will define new
H\"older spaces taking into account the highly non-isotropic
jumping geometry of the domains. We will give precise
H\"older estimates for solving $\overline\partial$ on the
domains in these spaces.}


\bigskip

\speaker{%
Xiaohua Zhu  (Australian National Univeristy, Australia)}
\title{Expending K\"ahler-Ricci solitons on  a class of quasi-projective
manifolds}

\abstract{%
In this talk I will discuss an existence problem  of complete
expending K\"ahler-Ricci solitons on a class of quasi-projective
manifolds. Our method is to reduce to solving certain complex
Monge-Amp\`ere equation on  a complete noncompact manifold. We will
fisrt construct an approximate metric of K\"ahler-Ricci soliton then we
study a family of complex Monge-Amp\`re equations with suitable Dirichlet boundary conditons. Our main result is analogous to  a theorem of Tian-Yau for the K\"ahler-Ricci flat metrics on quasi-projective manifolds.
}

\newpage

\speaker{%
Toshiki Mabuchi (Osaka University)}
\title{Stability, obstructions and extremal metrics}
\abstract{%
In a recent paper in JDG '02,  Donaldson estabished the
Chow-Mumford stability for a projective algebraic K\"ahler
manifold M of constant scalar curvature essentially
when the reductive part $R(G)$ of the maximal connected linear
algebraic subgroup $G$ of Aut$(M)$ is semisimple.

In this talk, I'll  give a generalization of his result to extremal
K\"ahler cases even when $R(G)$ has positive dimensional
center. In particular, a method of Donaldson
allows us to obtain the uniqueness, modulo the action of $G$,
of extremal K\"ahler metrics, if any, in a fixed integral
cohomology class.

If possible, I would like to talk also about a recent result of
Futaki on an obstruction to semistability of a polarized algebraic
manifold.
}


\bigskip

\speaker{%
Bo-Yong Chen (Nagoya University)}
\title{Recent results for the Bergman metric on K\"ahler manifolds}
\nopagebreak
\abstract{%
I will talk some recent results on the Bergman metric of complete K\"ahler manifolds. We will show the Bergman completeness of any hyperconvex manifold. An example which is Bergman complete but not holomorphically convex is given. We also give a localization principle for the Bergman kernel and metric of $C^2$ pseudoconvex domains in the complex projective space. 
}

\bigskip

\speaker{Mei-Chi Shaw (University of Notre Dame, USA)}
\title{Estimates for the $\bar\partial$-Neumann problem and nonexistence of Levi-flat hypersurfaces in $\mathbb{C}P^n$}
\abstract{%
Let $\Omega$ be a {\em  pseudoconvex } domain with $C^2$-smooth boundary
in $\mathbb{C}P^n$. We prove that the  $\bar\partial$-Neumann operator $N$ exists on
$\Omega$.
Furthermore, there exists an $\epsilon_0>0$ such that $N$ is regular on
$W^\epsilon ( {\Omega}) $ for
$\epsilon<\epsilon_0$.
We use the estimates to prove  the nonexistence of $C^{2, \alpha}$ real
 Levi-flat hypersurfaces  in $\mathbb{C}P^n$.
}

\end{document}





































 









 
